\label{der:3-1}
由列乘法定义，$A\vx=\sum_{j=1}^s x_j\va_j$。所以列组无关与$A\vx=0$只有零解是同一命题的两种写法。消元后若主元数$r<s$，令一个自由变量为1就产生非零解；若$r=s$，回代使全部未知量为零，故无关等价于满列秩。方形时满秩又等价于行列式非零，长方阵不使用这一判据。

若原组无关，加入$\vb$后有非零关系$\sum_i c_i\va_i+d\vb=0$。若$d=0$，原组无关迫使全部$c_i=0$，矛盾；故$d\ne0$，可解得$\vb=-\sum_i(\frac{c_i}{d})\va_i$。反向若$\vb=\sum_i b_i\va_i$，则$\sum_i b_i\va_i-\vb=0$是加入后的一条非零系数关系。若另有$\vb=\sum_i b_i'\va_i$，相减得$\sum_i(b_i-b_i')\va_i=0$，原组无关给$b_i=b_i'$，表示唯一。

最后，$n\times s$矩阵的主元最多$n$个；当$s>n$时不能满列秩，所以$n$维空间中超过$n$个向量必相关。
